% Introduction (Section 1) -- paste into your ML4H/PMLR Overleaf project.
% Requires natbib (the PMLR/JMLR style loads it) and amsmath for \alpha.
% Bib keys used: angelopoulos2024crc, min2023factscore, mohri2024conformal, cherian2024llm.

\section{Introduction}
\label{sec:intro}

Large language models (LLMs) are increasingly being used to support patients and
healthcare professionals by answering medical questions and assisting with
health-related information. Yet they state false claims as confidently as true
ones. A model might report that a drug should be taken at ``5000 mg per day'' just
as fluently as it reports that ``metformin treats type 2 diabetes,'' with no signal
that one of them is wrong. In most applications, factual errors are often minor and
carry limited consequences. In medicine, a wrong dose or a missed contraindication
can cause real harm. For these systems to be trusted in the clinic, it is not
enough to be usually right: we need a reliable bound on how often, and how
dangerously, they are wrong.

Raw confidence estimates from LLMs do not provide such a guarantee. Both softmax
probabilities and self-reported confidence scores are often poorly calibrated,
making them unreliable indicators of the true probability of error. Conformal
prediction offers a way to obtain this guarantee regardless of how well-calibrated
the underlying confidence score is. Given a held-out calibration set of examples
with known correctness, it converts any confidence score into a decision threshold
with a distribution-free guarantee. On new data, the error rate among kept
predictions stays below a target level $\alpha$ that the user chooses in advance.

Conformal risk control (CRC) extends this idea from prediction sets to any bounded
error measure that decreases as the threshold tightens \citep{angelopoulos2024crc},
such as the fraction of kept claims that are hallucinations; the target level
$\alpha$ then acts as an error budget the filter must respect. A medical answer is
rarely entirely true or entirely false, so it is first decomposed into atomic
claims, and a keep-or-flag decision is made for each claim on its own
\citep{min2023factscore}. Keeping only the claims confident enough to satisfy the
budget yields a conformal factuality guarantee \citep{mohri2024conformal}, and this
has recently been demonstrated on medical long-form question answering
\citep{cherian2024llm}. We take claim-level medical conformal factuality as our
starting point; it is the basis we build on, not our contribution.

Existing conformal factuality guarantees overlook a basic asymmetry: not all
hallucinations are equally harmful. A wrong dose or a missed contraindication is
dangerous, whereas an incorrect background fact, such as the year a drug was first
described, is benign. A single budget, however, treats every hallucination as an
identical unit of error. Because the guarantee constrains only the combined rate, a
filter can be highly accurate on benign claims and still have a very high error rate
on dangerous claims, as long as the benign claims, which are far more numerous, pull
the average back within budget. This is not hypothetical. Across two open models and
three medical QA benchmarks, a single global CRC threshold tuned to a 10 percent
marginal budget leaves the dangerous-claim hallucination rate between 0.068 and
0.094, well above a 0.05 dangerous target, on every model and dataset pair we test.

We address this by making the guarantee severity-aware. Rather than controlling one
harm-blind error rate, we partition claims into the two severity tiers above and
assign each tier its own conformal budget. Because the calibration procedure
enforces the budgets separately, holding the dangerous tier to a stricter target
than the benign tier directly forces a higher confidence threshold on dangerous
claims, rather than relying on severity to come out in the wash of an aggregate
rate. This reallocates strictness rather than raising it everywhere: the looser
benign budget lets the system surface more true benign information, even as the
stricter dangerous budget catches more dangerous hallucinations.

On the same two models and three datasets, severity-aware CRC brings the dangerous
hallucination rate to 0.049 or below, back within budget, at a modest cost in the
fraction of true claims retained. We report 95 percent confidence intervals from a
cluster bootstrap that resamples whole questions, since claims drawn from the same
answer are correlated. A score ablation further shows that the safety guarantee is a
property of the method rather than of any particular confidence score: even a random
score holds the dangerous tier within budget, and only the retention of true claims
suffers. Our contributions are:

\begin{itemize}
  \item We formulate \textbf{severity-aware (tiered) conformal risk control} for
  harm-stratified, claim-level factuality in medicine, separating the keep-or-flag
  confidence score from the per-tier harm budget.
  \item Across \textbf{two models and three medical QA datasets}, we show that a
  single global budget systematically exceeds the dangerous-claim target (0.068 to
  0.094) while severity-aware control restores it (0.049 or below), with
  cluster-bootstrap confidence intervals and per-split violation rates.
  \item We show that validity is \textbf{method-borne} through a score ablation, and
  that severity-awareness reallocates rather than uniformly increases strictness,
  keeping more true benign claims while catching more dangerous hallucinations.
\end{itemize}
